\documentclass[11pt]{amsart}
\usepackage[T1]{fontenc}
\usepackage{lmodern,microtype,amsmath,amssymb,amsthm,mathtools,enumitem,booktabs,longtable}
\usepackage[margin=1.1in]{geometry}
\usepackage[hidelinks]{hyperref}
\usepackage{xurl}

\newtheorem{theorem}{Theorem}[section]
\newtheorem{lemma}[theorem]{Lemma}
\newtheorem{proposition}[theorem]{Proposition}
\newtheorem{corollary}[theorem]{Corollary}
\theoremstyle{definition}
\newtheorem{definition}[theorem]{Definition}
\theoremstyle{remark}
\newtheorem{remark}[theorem]{Remark}
\numberwithin{equation}{section}

\newcommand{\RR}{\mathbb{R}}
\newcommand{\cF}{\mathcal{F}}
\newcommand{\cP}{\mathbb{P}}
\newcommand{\fat}{\operatorname{fat}}
\newcommand{\sgn}{\operatorname{sgn}}

\hypersetup{pdftitle={The fat-shattering dimension of norm-bounded two-layer ReLU networks in fixed dimension},pdfauthor={Yitzchak Shmalo}}

\title[Fat-shattering of two-layer ReLU networks in fixed dimension]{The fat-shattering dimension of norm-bounded two-layer ReLU networks in fixed dimension}
\author{Yitzchak Shmalo}
\address{Einstein Institute of Mathematics, The Hebrew University of Jerusalem, Jerusalem, Israel}
\email{yitzchak.shmalo@gmail.com}
\date{November 2026}
\subjclass[2020]{68Q32, 41A46, 68T07, 52A21}
\keywords{fat-shattering dimension, sample complexity, ReLU networks, variation spaces, zonoids}

\begin{document}
\begin{abstract}
Consider the functions $x\mapsto u^{T}\sigma(Wx)$ on the ball of radius $b_x$ in $\RR^d$, where $\sigma$ is the ReLU function, $W$ has $n$ rows, $\|W\|_F\le B$ and $|u|\le b$, and put $R=bBb_x/\varepsilon$. For every width and dimension the fat-shattering dimension of this class at scale $\varepsilon$ is at most $R^2$ up to logarithmic factors, and Vardi, Shamir and Srebro asked whether a better bound holds when $d$ is small. We prove that for fixed $d\ge2$ the fat-shattering dimension is at most a constant times $R^{2(d-1)/(d+2)}\log(2+R)$ for every width, and at least a constant times $R^{2(d-1)/(d+2)}$ when $R$ exceeds a constant depending on $d$ and the width is at least that large. The upper bound writes the networks through support functions of zonoids, as Bach did, and uses the approximation of zonoids by zonotopes due to Matou\v{s}ek and Siegel. The lower bound restricts the networks to an affine hyperplane and uses the embedding of Sobolev spaces into the variation space of ReLU networks with biases due to Mao, Siegel and Xu.
\end{abstract}
\maketitle

\section{Introduction}\label{sec:intro}

Fix $d\ge2$, a width $n\ge1$ and constants $b,B,b_x>0$. Let $\sigma(z)=\max\{z,0\}$, applied coordinatewise, let $X=\{x\in\RR^d:|x|\le b_x\}$, and let
\[
\cF_n=\bigl\{x\mapsto u^{T}\sigma(Wx)\;:\;u\in\RR^n,\ W\in\RR^{n\times d},\ |u|\le b,\ \|W\|_F\le B\bigr\},
\]
a class of functions on $X$; here $|\cdot|$ is the Euclidean norm and $\|\cdot\|_F$ the Frobenius norm. The sample complexity of learning with such a class is governed by its fat-shattering dimension, introduced by Kearns and Schapire~\cite{KS}; its role in learning was established by Alon, Ben-David, Cesa-Bianchi and Haussler~\cite{ABCH} and by Bartlett~\cite{Bartlett98}.

\begin{definition}
Let $\cF$ be a class of real functions on a set $X$ and $\varepsilon>0$. Points $x_1,\dots,x_N\in X$ are \emph{$\varepsilon$-shattered} by $\cF$ if there are $r_1,\dots,r_N\in\RR$, the levels, such that for every $\varsigma\in\{-1,1\}^N$ some $f\in\cF$ satisfies $\varsigma_j(f(x_j)-r_j)\ge\varepsilon$ for $j=1,\dots,N$. The fat-shattering dimension $\fat_\varepsilon(\cF)$ is the largest such $N$, or $+\infty$ if there is no largest one.
\end{definition}

Write $R=bBb_x/\varepsilon$. Neyshabur, Tomioka and Srebro~\cite{NTS} and Golowich, Rakhlin and Shamir~\cite{GRS} bound the Rademacher complexity of $\cF_n$, Vardi, Shamir and Srebro~\cite[Theorem~2]{VSS} extend the bound to Lipschitz activations, and Bartlett, Foster and Telgarsky~\cite{BFT} prove bounds for deep networks in terms of spectral norms. For every width $n$ and every dimension $d$ the Rademacher bound gives $\fat_\varepsilon(\cF_n)\le CR^2$ up to logarithmic factors. Vardi, Shamir and Srebro also prove lower bounds that grow linearly in $R$: for the activation $z\mapsto\sigma(z-\beta)$ with a suitable $\beta$ when $d$ exceeds a universal constant~\cite[Theorem~3]{VSS}, and for $\sigma$ itself when $d$ exceeds a factor polylogarithmic in the parameters~\cite[Corollary~1]{VSS}. They leave open whether a bound of smaller order than $R^2$ holds when $d$ is fixed~\cite[Section~3.2]{VSS}.

Two results replace the power $2$ of $R$ by a smaller one when $d$ is fixed. Liu, Dadi and Cevher~\cite[Proposition~5]{LDC} bound, for $d>5$, the logarithm of the $\ell_2$ empirical covering number at scale $\eta$ of the two-layer ReLU networks $\frac1m\sum_{k\le m}a_k\sigma(w_k\cdot x)$ on the cube $\{\|x\|_\infty\le1\}$ with $\frac1m\sum_k|a_k|\,\|w_k\|_1\le1$ by a constant times $\eta^{-2d/(d+2)}$. Since a set shattered at scale $\varepsilon$ carries an $\varepsilon$-separated family of size exponential in its cardinality, this gives $\fat_\varepsilon(\cF_n)\le C_dR^{2d/(d+2)}$. The uniform approximation theorem of Siegel~\cite[Theorem~2]{Siegel} gives, for every $d\ge2$ and every $\beta\ge0$, the bound $\fat_\varepsilon(\cF_n^{\beta})\le C_dR^{2d/(d+3)}\log(2+R)$ (Proposition~\ref{prop:bias}), where $\cF_n^{\beta}$ consists of the networks $x\mapsto u^{T}\sigma(Wx-\beta\mathbf 1)$ with a common bias $\beta$ and the same constraints on $u$ and $W$. The case $\beta=0$ is $\cF_n$. Both arguments treat $\sigma(w\cdot x)$ as a ridge function of $d$ variables. Neither uses that it is positively homogeneous. The lower bound of~\cite[Theorem~3]{VSS} holds for every $d$ above a universal constant but concerns the activation $z\mapsto\sigma(z-\beta)$, and the one for $\sigma$ itself~\cite[Corollary~1]{VSS} requires $d$ to exceed a factor polylogarithmic in the parameters. The Sobolev embedding of Mao, Siegel and Xu~\cite{MSX} produces functions of small variation norm (Section~\ref{sec:prelim}), not networks of a given width with $|u|\le b$ and $\|W\|_F\le B$.

A neuron $\sigma(w\cdot x)$ of $\cF_n$ is positively homogeneous, and on the affine hyperplane $\{x_d=b_x/\sqrt2\}$ it is a neuron $\sigma(w'\cdot y+t)$ in the remaining $e=d-1$ variables. Theorem~\ref{thm:main} shows that for fixed $d$ the power of $R$ in $\fat_\varepsilon(\cF_n)$ is $2(d-1)/(d+2)$: the upper bound $C_dR^{2(d-1)/(d+2)}\log(2+R)$ holds for every width, and the lower bound $c_dR^{2(d-1)/(d+2)}$ holds for $R\ge C'_d$ and every width $n\ge c_dR^{2(d-1)/(d+2)}$. The power $2(d-1)/(d+2)$ equals $2e/(e+3)$, the power in the bound for $\cF_n^{\beta}$ with $e$ in place of $d$. The upper bound combines known approximation theorems; the lower bound is new.

\begin{theorem}\label{thm:main}
For every $d\ge2$ there are constants $c_d,C_d,C'_d>0$, depending only on $d$, such that the following hold for all $n\ge1$ and $b,B,b_x,\varepsilon>0$, with $R=bBb_x/\varepsilon$.
\begin{enumerate}[label=(\alph*)]
\item If $R<1$ then $\fat_\varepsilon(\cF_n)=0$, and for all $R$,
\[
\fat_\varepsilon(\cF_n)\le C_d\,R^{2(d-1)/(d+2)}\log(2+R).
\]
\item If $R\ge C'_d$ then
\[
\fat_\varepsilon(\cF_n)\ge\min\bigl\{n,\ \lfloor c_d\,R^{2(d-1)/(d+2)}\rfloor\bigr\}.
\]
\end{enumerate}
\end{theorem}

The power $2(d-1)/(d+2)$ equals $1/2$ for $d=2$ and $1$ for $d=4$, and it increases to $2$, the power of $R$ in the dimension-free bound, as $d\to\infty$. For $R\ge C'_d$ and width at least $c_dR^{2(d-1)/(d+2)}$ the two bounds agree up to the logarithm. Let $\cF_n^{\mathrm{sp}}$ be the class defined as $\cF_n$ but with the spectral norm $\|W\|\le B$ in place of $\|W\|_F\le B$. Vardi, Shamir and Srebro~\cite[Theorem~1]{VSS} proved that under spectral-norm control the fat-shattering dimension grows with the width once $d$ is at least a polynomial in $b$, $B$, $b_x$, $n$ and $1/\varepsilon$; the following corollary, proved at the end of Section~\ref{sec:upper}, shows that for fixed $d$ it does not.

\begin{corollary}\label{cor:spectral}
Theorem~\ref{thm:main} holds for $\cF_n^{\mathrm{sp}}$, with $C_d$ replaced by a larger constant depending only on $d$.
\end{corollary}

By positive homogeneity every $f\in\cF_n$ is a combination $\sum_ic_i\sigma(\omega_i\cdot x)$ of unit-norm neurons with $\sum_i|c_i|\le bB$. Since $\sigma(z)=(|z|+z)/2$, $f$ is, up to a linear function, a difference of two support functions of zonoids in $\RR^d$ whose generating measures have total mass at most $bB$; this connection between ReLU networks and zonoids is due to Bach~\cite{Bach}. Zonoids in $\RR^d$ are approximated by zonotopes with $m$ summands with error of order $m^{-1/2-3/(2(d-1))}$ in the uniform norm on the sphere. For $d\ge5$ this is due to Matou\v{s}ek~\cite{Matousek}, and for $d=3,4$, where Matou\v{s}ek's bound carries a factor $\sqrt{\log m}$, to Siegel~\cite[Theorem~1]{Siegel}; the rate is optimal by Bourgain, Lindenstrauss and Milman~\cite{BLM}. For $d=2$ it is the classical approximation of a planar symmetric convex body by polygons (McClure and Vitale~\cite{MV}; Lemma~\ref{lem:planar}). This is the rate of~\cite[Theorem~2]{Siegel} for combinations of the functions in $\cP^{d-1}$ (Section~\ref{sec:prelim}); Siegel remarks after his Theorem~3 that his Theorem~1 is essentially the case $k=1$ of his Theorem~2. It gives a uniform cover of $\cF_n$ whose logarithm is of order $R^{2(d-1)/(d+2)}\log R$, and a uniform cover at scale below $\varepsilon$ bounds $\fat_\varepsilon$ by the logarithm of its size.

The lower bound uses points on the affine hyperplane $\{x_d=b_x/\sqrt2\}$. There $\sigma(w\cdot x)$ restricts to a neuron with a bias in dimension $d-1$, so on these points $\cF_n$ contains the combinations of the functions in $\cP^{d-1}$ whose variation norm, in the sense of Bach~\cite{Bach} and Siegel and Xu~\cite{SX} (Section~\ref{sec:prelim}), is at most of order $bBb_x$. Mao, Siegel and Xu~\cite{MSX} proved that the Sobolev space of order $s=(e+3)/2$ on the unit ball of $\RR^{e}$ embeds into the variation space $\mathcal K_1(\cP^e)$. A sum of $N$ bumps of height $2\varepsilon$ and width $\delta$, with arbitrary signs, has Sobolev norm of order $\varepsilon\delta^{-s}$, so its variation norm is below $bBb_x$ once $\delta$ is of order $R^{-1/s}$, and the bumps then number of order $R^{(d-1)/s}=R^{2(d-1)/(d+2)}$. A Carath\'eodory argument (Lemma~\ref{lem:finite}) represents the values at the $N$ points by a network of width at most $N$ with $|u|\le b$ and $\|W\|_F\le B$.

For $\cF_n^{\beta}$, the class of~\cite[Theorem~3]{VSS}, the approximation theorem of Siegel for ReLU networks with biases~\cite[Theorem~2]{Siegel} gives the following bound, proved in Section~\ref{sec:bias}.

\begin{proposition}\label{prop:bias}
For every $d\ge2$ there is $C_d$ such that for all $\beta\ge0$, $n\ge1$ and $b,B,b_x,\varepsilon>0$, $\fat_\varepsilon(\cF_n^{\beta})=0$ if $R<1$ and $\fat_\varepsilon(\cF_n^{\beta})\le C_dR^{2d/(d+3)}\log(2+R)$ for all $R$.
\end{proposition}

The power $2d/(d+3)$ is larger than $2(d-1)/(d+2)$. From below, Theorem~\ref{thm:main}(b) is the case $\beta=0$, and Vardi, Shamir and Srebro~\cite[Theorem~3]{VSS} prove a lower bound linear in $R$ for a suitable $\beta$ when $d$ exceeds a universal constant. The construction of Section~\ref{sec:lower} does not transfer to $\cF_n^{\beta}$ when $\beta>0$: a neuron $\sigma(w\cdot x-\beta)$ that is not zero on $X$ has $|w|>\beta/b_x$, so fewer than $(Bb_x/\beta)^2$ neurons of a network in $\cF_n^{\beta}$ are not zero on $X$. From above, the same observation places $\cF_n^{\beta}$ inside $\cF_n$ in dimension $d+1$: replacing each row $w$ with $|w|>\beta/b_x$ by $(w,-\beta/b_x)$, each other row by $0$, and $x$ by $(x,b_x)$ multiplies $\|W\|_F$ and $b_x$ by at most $\sqrt2$, so Theorem~\ref{thm:main}(a) in dimension $d+1$ gives Proposition~\ref{prop:bias} again, with $2R$ in place of $R$. The infimum of the numbers $a$ for which $\sup_{\beta\ge0}\fat_\varepsilon(\cF_n^{\beta})\le C_dR^{a}\log(2+R)$ holds for all $R$ and $n$ lies between $2(d-1)/(d+2)$ and $2d/(d+3)$; its value is open.

The metric entropy of the variation space has been studied in $L_2$ and in the uniform norm. Siegel and Xu~\cite{SX} determined the $L_2$ metric entropy of the unit ball of $\mathcal K_1(\cP^d)$, of order $\eta^{-2d/(d+3)}$, and Siegel~\cite[Theorems~2 and~3]{Siegel} proved the matching approximation rate in the uniform norm, which gives a uniform entropy bound of the same order up to a logarithm (Section~\ref{sec:bias}); earlier uniform rates are due to Makovoz~\cite{Makovoz} and Bach~\cite{Bach}. Parhi and Nowak~\cite{PN} used such entropy bounds to obtain near-minimax rates for estimation with shallow ReLU networks, and Liang, Qiao, Wang and Parhi~\cite{LQWP} constructed packings of shallow ReLU networks that concentrate near the boundary of the domain to prove lower bounds for flat minima.

\section{Preliminaries}\label{sec:prelim}

For $e\ge1$ let $B^e$ be the closed unit ball of $\RR^e$ and $S^{e-1}$ its boundary sphere, with $S^0=\{-1,1\}$. Following~\cite{SX,MSX}, let
\[
\cP^e=\{\sigma(\omega\cdot y+t)\;:\;\omega\in S^{e-1},\ t\in[-1,1]\},
\]
regarded as a set of functions on $B^e$. Let $B_1(\cP^e)$ be the closure in $L_2(B^e)$ of the set of finite combinations $\sum_ia_ig_i$ with $g_i\in\cP^e$ and $\sum_i|a_i|\le1$, and for $g\in L_2(B^e)$ let $\|g\|_{\mathcal K_1(\cP^e)}=\inf\{c>0:g\in cB_1(\cP^e)\}$, the variation norm of $g$; the variation space of ReLU networks with biases, $\mathcal K_1(\cP^e)$, is the set of $g$ for which it is finite. Every element of $\cP^e$ is $1$-Lipschitz on $B^e$ and bounded by $2$ in absolute value. For $s\ge0$ and $g\in L_2(\RR^e)$ with Fourier transform $\hat g$ let
\[
\|g\|_{H^s}^2=\int_{\RR^e}(1+|\xi|)^{2s}|\hat g(\xi)|^2\,d\xi .
\]
This norm is equivalent to the norm of the space denoted $W^s(L_2(\RR^e))$ in~\cite{MSX}. The next two facts are standard; the reduction in Lemma~\ref{lem:reduce} goes back to Neyshabur, Tomioka and Srebro~\cite{NTS}.

\begin{lemma}\label{lem:cover}
Let $\cF$ be a class of real functions on a set $X$, let $\varepsilon>0$, and let $G$ be a finite set of real functions on $X$ such that for every $f\in\cF$ there is $g\in G$ with $\sup_X|f-g|<\varepsilon$. Then $\fat_\varepsilon(\cF)\le\log_2|G|$.
\end{lemma}

\begin{proof}
Let $x_1,\dots,x_N$ be $\varepsilon$-shattered with levels $r_1,\dots,r_N$. For $\varsigma\in\{-1,1\}^N$ choose $f_\varsigma\in\cF$ with $\varsigma_j(f_\varsigma(x_j)-r_j)\ge\varepsilon$ for all $j$, and $g_\varsigma\in G$ with $\sup_X|f_\varsigma-g_\varsigma|<\varepsilon$. Then $\varsigma_j(g_\varsigma(x_j)-r_j)>0$ for all $j$, so $\varsigma$ is the vector of signs of $(g_\varsigma(x_j)-r_j)_j$, and $\varsigma\mapsto g_\varsigma$ is injective. Hence $2^N\le|G|$.
\end{proof}

\begin{lemma}\label{lem:reduce}
Every $f\in\cF_n$ can be written as $f(x)=\sum_{i=1}^{n}c_i\sigma(\omega_i\cdot x)$ with $\omega_i\in S^{d-1}$ and $\sum_i|c_i|\le bB$. In particular $|f(x)|\le bB|x|$, and $\fat_\varepsilon(\cF_n)=0$ when $R<1$.
\end{lemma}

\begin{proof}
Let $f(x)=\sum_iu_i\sigma(w_i\cdot x)$, where $w_i$ is the $i$th row of $W$. Rows with $w_i=0$ contribute nothing; for the others put $\omega_i=w_i/|w_i|$ and $c_i=u_i|w_i|$, so that $u_i\sigma(w_i\cdot x)=c_i\sigma(\omega_i\cdot x)$ by positive homogeneity. By the Cauchy--Schwarz inequality $\sum_i|c_i|\le|u|\,\|W\|_F\le bB$, and $|\sigma(\omega_i\cdot x)|\le|x|$. If a single point $x\in X$ is $\varepsilon$-shattered with level $r$, there are $f_+,f_-\in\cF_n$ with $f_+(x)\ge r+\varepsilon$ and $f_-(x)\le r-\varepsilon$, so $2\varepsilon\le2bBb_x$, that is $R\ge1$.
\end{proof}

\section{The upper bound}\label{sec:upper}

For a finite positive measure $\tau$ on $S^{d-1}$ let $h_\tau(x)=\int_{S^{d-1}}|x\cdot y|\,d\tau(y)$; this is the support function of a zonoid, and every zonoid centred at the origin arises in this way. A measure with $m$ atoms gives the support function of a zonotope with $m$ summands, and the approximation of $h_\tau$ by such measures is the following statement. For $d\ge5$ it is due to Matou\v{s}ek~\cite{Matousek}, for $d=3,4$ to Siegel~\cite[Theorem~1]{Siegel}, and for $d=2$ it is classical~\cite{MV}; we include a proof for $d=2$ with explicit constants.

\begin{lemma}\label{lem:zonoid}
For every $d\ge2$ there is $A_d$ such that for every probability measure $\tau$ on $S^{d-1}$ and every $m\ge1$ there is a positive measure $\tau'$ on $S^{d-1}$ with at most $m$ atoms and total mass at most $1$ such that
\[
\sup_{x\in S^{d-1}}|h_\tau(x)-h_{\tau'}(x)|\le A_d\,m^{-(d+2)/(2(d-1))}.
\]
\end{lemma}

\begin{proof}
For $d\ge3$ this is~\cite[Theorem~1]{Siegel}, which for $d\ge5$ recovers the bound of~\cite{Matousek}, stated there for the sphere $S^{e}\subset\RR^{e+1}$ with the exponent $1/2+3/(2e)$; with $e=d-1\ge2$ the exponent is $(d+2)/(2(d-1))$, and the measure provided there is a probability measure. For $d=2$ the exponent is $2$, and the statement is Lemma~\ref{lem:planar} below.
\end{proof}

\begin{lemma}\label{lem:planar}
Let $\tau$ be a probability measure on $S^{1}$ and $m\ge1$. There are $v_1,\dots,v_m\in\RR^2$ with $\sum_i|v_i|\le1$ and $\sup_{x\in S^1}|h_\tau(x)-\sum_{i=1}^m|x\cdot v_i||\le40\,m^{-2}$.
\end{lemma}

\begin{proof}
The function $h=h_\tau$ is even, positively homogeneous and subadditive, so it is the support function of a compact convex set $Z\subset\RR^2$ symmetric about the origin, and $h(x)\le|x|$ gives $Z\subset B^2$. If $Z$ is a segment $[-v,v]$ then $h(x)=|x\cdot v|$ and $|v|\le1$, so assume that $Z$ has nonempty interior. Let $\Lambda$ denote the perimeter of $Z$. By Cauchy's formula $\Lambda=\int_0^{2\pi}h(\cos\theta,\sin\theta)\,d\theta$, and since $\int_0^{2\pi}|\cos(\theta-\psi)|\,d\theta=4$ for every $\psi$, $\Lambda=4\tau(S^1)=4$. Since $0\le h\le1$ on $S^1$, the choice $v_1=\dots=v_m=0$ proves the bound for $m\le4$, so let $m\ge5$.

Fix $p_0\in\partial Z$ and let $\Gamma$ be the arc of $\partial Z$ from $p_0$ to $-p_0$ in the counterclockwise direction. Along $\Gamma$ the arc length increases by $\Lambda/2=2$ and the direction of the tangent turns by $\pi$ in total, counting the turns at corners. Let $\mu$ be the sum of the arc-length measure on $\Gamma$ and the measure of turning (the distributional derivative of the angle of the right tangent), a measure of total mass at most $2\pi$. Taking quantiles of $\mu$, choose points $p_0,p_1,\dots,p_m=-p_0$ in order along $\Gamma$, not necessarily distinct, such that $\mu$ gives mass at most $2\pi/m$ to every open arc $(p_{k-1},p_k)$. Put $p_{m+k}=-p_k$ and let $Z_m$ be the convex hull of $p_0,\dots,p_{2m-1}$, a polygon symmetric about the origin with $Z_m\subset Z$.

Consider an open arc between consecutive points $p=p_{k-1}$ and $p'=p_k$, of length $\ell\le2\pi/m$, along which the tangent turns by $\alpha\le2\pi/m<\pi/2$. The unit tangents along the arc lie in an angular interval of length $\alpha$, and the vectors $q-p$ and $p'-p$, for $q$ on the arc, are integrals of these tangents, so the angle between them is at most $\alpha$ and both make an acute angle with $p'-p$; similarly for $p'-q$. Hence the orthogonal projection of $q$ onto the line through $p$ and $p'$ lies on the segment $[p,p']$, and the distance from $q$ to that segment is at most $|q-p|\sin\alpha\le\ell\alpha\le4\pi^2m^{-2}$. Every point of $\partial Z$ is therefore within $4\pi^2m^{-2}$ of $Z_m$, and since $Z_m\subset Z$ we get $0\le h-h_{Z_m}\le4\pi^2m^{-2}<40\,m^{-2}$ on $S^1$, where $h_{Z_m}$ is the support function of $Z_m$.

A convex polygon symmetric about the origin has edge vectors $\zeta_1,\dots,\zeta_{m'},-\zeta_1,\dots,-\zeta_{m'}$ in cyclic order, with $m'\le m$ here, and it is the Minkowski sum of the segments $[-\zeta_i/2,\zeta_i/2]$. Hence $h_{Z_m}(x)=\sum_{i\le m'}|x\cdot\zeta_i|/2$. Put $v_i=\zeta_i/2$ for $i\le m'$ and $v_i=0$ otherwise. Then $\sum_i|v_i|$ is a quarter of the perimeter of $Z_m$, which is at most $\Lambda$, so $\sum_i|v_i|\le\Lambda/4=1$.
\end{proof}

\begin{lemma}\label{lem:Hcover}
For every $d\ge2$ there is $C_d$ such that the following holds. Let $\mathcal H$ be the set of functions on $B^d$ of the form $h(x)=\sum_ic_i\sigma(\omega_i\cdot x)$, a finite sum with $\omega_i\in S^{d-1}$ and $\sum_i|c_i|\le1$. For every $\eta\in(0,1]$ there is a finite set $G$ of functions on $B^d$ with $\log|G|\le C_d\,\eta^{-2(d-1)/(d+2)}\log(2/\eta)$ such that every $h\in\mathcal H$ satisfies $\sup_{B^d}|h-g|\le\eta$ for some $g\in G$.
\end{lemma}

\begin{proof}
Write $\gamma=(d+2)/(2(d-1))$. Let $h\in\mathcal H$. Since $\sigma(z)=(|z|+z)/2$,
\[
h(x)=\tfrac12\bigl(P\,h_{\tau_+}(x)-Q\,h_{\tau_-}(x)\bigr)+\mathbf m\cdot x,\qquad \mathbf m=\tfrac12\sum_ic_i\omega_i,
\]
where $P=\sum_{c_i>0}c_i$, $Q=\sum_{c_i<0}|c_i|$, $P+Q\le1$, $|\mathbf m|\le1/2$, and $\tau_\pm$ are probability measures on $S^{d-1}$ (if $P=0$ or $Q=0$ the corresponding term is absent). Let $m$ be the least integer with $A_dm^{-\gamma}\le\eta/2$. By Lemma~\ref{lem:zonoid} there are positive measures $\tau'_\pm$ with at most $m$ atoms and mass at most $1$ such that $|h_{\tau_\pm}-h_{\tau'_\pm}|\le\eta/2$ on $S^{d-1}$, and by homogeneity $|h_{\tau_\pm}(x)-h_{\tau'_\pm}(x)|\le\eta|x|/2$ on $B^d$. Hence
\[
\Bigl|h(x)-\sum_{i=1}^{2m}a_i|v_i\cdot x|-\mathbf m\cdot x\Bigr|\le\frac\eta4\qquad(x\in B^d)
\]
for some $v_i\in S^{d-1}$ and $a_i\in\RR$ with $\sum_i|a_i|\le1$.

Let $\Delta=\eta/32$. Let $N_1$ be a $\Delta$-net of $S^{d-1}$ and $N_2$ a $\Delta$-net of the ball of radius $1/2$ in $\RR^d$, each of cardinality at most $(3/\Delta)^d$, and let $N_3=(\Delta/(4m))\mathbb Z\cap[-2,2]$. Let $G$ be the set of functions $x\mapsto\sum_{i=1}^{2m}a'_i|v'_i\cdot x|+\mathbf m'\cdot x$ on $B^d$ with $a'_i\in N_3$, $v'_i\in N_1$ and $\mathbf m'\in N_2$. Replacing $v_i$, $a_i$ and $\mathbf m$ by nearest points of the nets changes the function on $B^d$ by at most $\sum_i|a_i|\Delta+2m\cdot\Delta/(8m)+\Delta\le4\Delta$, since $||v\cdot x|-|v'\cdot x||\le|v-v'|$ and $|v'\cdot x|\le1$ for $|x|\le1$. As $4\Delta<\eta/4$, every $h\in\mathcal H$ is within $\eta/2$ of an element of $G$. Finally
\[
\log|G|\le2m\bigl(d\log(3/\Delta)+\log(16m/\Delta+1)\bigr)+d\log(3/\Delta),
\]
and $m\le(2A_d/\eta)^{1/\gamma}+1$ with $1/\gamma=2(d-1)/(d+2)$, which gives the stated bound since $\log m$ and $\log(1/\Delta)$ are at most constants times $\log(2/\eta)$.
\end{proof}

\begin{proof}[Proof of Theorem~\ref{thm:main}(a)]
The case $R<1$ is Lemma~\ref{lem:reduce}. Let $R\ge1$ and $f\in\cF_n$. By Lemma~\ref{lem:reduce}, $h(y)=f(b_xy)/(bBb_x)$ belongs to $\mathcal H$. Let $\eta=1/(2R)\le1/2$ and let $G$ be the set of Lemma~\ref{lem:Hcover}. For $g\in G$ with $\sup_{B^d}|h-g|\le\eta$, the function $x\mapsto bBb_x\,g(x/b_x)$ is within $bBb_x\eta=\varepsilon/2<\varepsilon$ of $f$ on $X$. Lemma~\ref{lem:cover} applied to these rescaled functions gives
\[
\fat_\varepsilon(\cF_n)\le\log_2|G|\le C_d(2R)^{2(d-1)/(d+2)}\log(4R)/\log2,
\]
which is at most a constant depending on $d$ times $R^{2(d-1)/(d+2)}\log(2+R)$.
\end{proof}

\begin{proof}[Proof of Corollary~\ref{cor:spectral}]
Since $\|W\|\le\|W\|_F$ and $\|W\|_F\le\sqrt d\,\|W\|$ for $W\in\RR^{n\times d}$, the class $\cF_n$ is contained in $\cF_n^{\mathrm{sp}}$, and $\cF_n^{\mathrm{sp}}$ is contained in the class $\cF_n$ defined with $\sqrt d\,B$ in place of $B$. The lower bound for $\cF_n$ therefore holds for $\cF_n^{\mathrm{sp}}$. For $R\ge1$ the upper bound for $\cF_n^{\mathrm{sp}}$ follows from part~(a) applied with $\sqrt d\,R$ in place of $R$. For $R<1$, if $\|W\|\le B$ and $|u|\le b$ then $|u^{T}\sigma(Wx)|\le|u|\,|Wx|\le bB|x|$, and the last argument in the proof of Lemma~\ref{lem:reduce} gives $\fat_\varepsilon(\cF_n^{\mathrm{sp}})=0$.
\end{proof}

\section{The lower bound}\label{sec:lower}

Throughout this section $e=d-1\ge1$ and $s=(e+3)/2=(d+2)/2$.

\begin{lemma}[Mao, Siegel and Xu~\cite{MSX}]\label{lem:embed}
There is $C_e$ such that $\|g|_{B^e}\|_{\mathcal K_1(\cP^e)}\le C_e\|g\|_{H^s}$ for every $g\in C_c^\infty(\RR^e)$.
\end{lemma}

\begin{proof}
By~\cite[Theorem~1]{MSX} with $k=1$, the space $W^s(L_2(B^e))$ embeds continuously into $\mathcal K_1(\cP^e)$; in~\cite{MSX} the norm of $W^s(L_2(B^e))$ is defined as the infimum of the norms in $W^s(L_2(\RR^e))$ of the extensions, and the latter norm is bounded by a constant times $\|g\|_{H^s}$.
\end{proof}

\begin{lemma}\label{lem:bumps}
Let $\varphi\in C_c^\infty(\RR^e)$ be supported in the open unit ball. There is $C_{e,\varphi}$ such that for $\delta\in(0,1]$, points $y_1,\dots,y_N\in\RR^e$ with $|y_j-y_l|\ge2\delta$ for $j\ne l$ and $N\delta^e\le1$, and signs $\varsigma_j\in\{-1,1\}$, the function $g=\sum_j\varsigma_j\varphi((\cdot-y_j)/\delta)$ satisfies $\|g\|_{H^s}\le C_{e,\varphi}\delta^{-s}$.
\end{lemma}

\begin{proof}
The supports of the summands are disjoint. For an integer $k\ge0$ we have $(1+|\xi|)^{2k}\le C\sum_{|\alpha|\le k}|\xi^\alpha|^2$ with $C$ depending on $e$ and $k$, so by Plancherel's theorem $\|g\|_{H^k}^2$ is at most a constant times $\sum_{|\alpha|\le k}\|\partial^\alpha g\|_{L_2}^2$. By disjointness and scaling, $\|\partial^\alpha g\|_{L_2}^2=N\delta^{e-2|\alpha|}\|\partial^\alpha\varphi\|_{L_2}^2\le\delta^{-2k}\|\partial^\alpha\varphi\|_{L_2}^2$ for $|\alpha|\le k$. Hence $\|g\|_{H^k}\le C\delta^{-k}$. If $s=k+\theta$ with $\theta\in(0,1)$, H\"older's inequality gives $\|g\|_{H^s}\le\|g\|_{H^k}^{1-\theta}\|g\|_{H^{k+1}}^{\theta}\le C\delta^{-s}$.
\end{proof}

\begin{lemma}\label{lem:finite}
Let $g$ be a continuous function on $B^e$ with $\|g\|_{\mathcal K_1(\cP^e)}\le M$, and let $y_1,\dots,y_N\in B^e$. There are $N'\le N$ elements $\sigma(\omega_i\cdot y+t_i)$ of $\cP^e$ and coefficients $a_i$ with $\sum_i|a_i|\le M$ such that $g(y_j)=\sum_{i=1}^{N'}a_i\sigma(\omega_i\cdot y_j+t_i)$ for $j=1,\dots,N$.
\end{lemma}

\begin{proof}
Let $K\subset\RR^N$ be the set of vectors $(\sigma(\omega\cdot y_j+t))_{j\le N}$ and $-(\sigma(\omega\cdot y_j+t))_{j\le N}$ with $\omega\in S^{e-1}$ and $t\in[-1,1]$. It is compact, so its convex hull $\operatorname{conv}K$ is compact. If $M=0$ then $g=0$ in $L_2(B^e)$, hence everywhere, and there is nothing to prove; let $M>0$. By the definition of the norm, $g\in M'B_1(\cP^e)$ for every $M'>M$, so there are finite combinations $g_k$ with coefficient sums at most $M'$ and $g_k\to g$ in $L_2(B^e)$. The $g_k$ are $M'$-Lipschitz and bounded by $2M'$, so by the Arzel\`a--Ascoli theorem a subsequence converges uniformly on $B^e$ to a continuous function, which equals $g$ almost everywhere and therefore everywhere on $B^e$. Hence $v=(g(y_j))_j$ is a limit of points of $M'\operatorname{conv}K$, so $v\in M'\operatorname{conv}K$ for every $M'>M$, and by compactness $v\in M\operatorname{conv}K$. If $v=0$ take $N'=0$. Otherwise let $\lambda\in(0,M]$ be the least number with $v\in\lambda\operatorname{conv}K$; then $v/\lambda$ lies on the boundary of $\operatorname{conv}K$, in a face of dimension at most $N-1$, and by Carath\'eodory's theorem it is a convex combination of at most $N$ points of $K$. Multiplying by $\lambda$ gives the claim.
\end{proof}

\begin{proof}[Proof of Theorem~\ref{thm:main}(b)]
Fix $\varphi\in C_c^\infty(\RR^e)$ supported in the open unit ball with $\varphi(0)=1$, and let $C_1=C_eC_{e,\varphi}$ with the constants of Lemmas~\ref{lem:embed} and~\ref{lem:bumps}. Put $\delta=(4C_1/R)^{1/s}$ and $C'_d=8^s\cdot4C_1$, so that $\delta\le1/8$ when $R\ge C'_d$. The points of the lattice $2\delta\mathbb Z^e$ in the ball of radius $1/2$ number at least $c\,\delta^{-e}$ for a constant $c=c(e)>0$, they are at distance at least $2\delta$ from each other, and since the balls of radius $\delta$ around them are disjoint and lie in $B^e$, any $N$ of them satisfy $N\delta^e\le1$. Let $c_d=c\,(4C_1)^{-e/s}$ and $N=\min\{n,\lfloor c_dR^{e/s}\rfloor\}$, choose $N$ of these points $y_1,\dots,y_N$, and note $e/s=2(d-1)/(d+2)$. Put
\[
x_j=\frac{b_x}{\sqrt2}(y_j,1)\in\RR^d,\qquad j=1,\dots,N,
\]
so that $|x_j|\le b_x$. We show that $x_1,\dots,x_N$ are $\varepsilon$-shattered by $\cF_n$ with all levels $r_j=0$.

Let $\varsigma\in\{-1,1\}^N$ and $g=2\varepsilon\sum_j\varsigma_j\varphi((\cdot-y_j)/\delta)$. Then $g(y_j)=2\varepsilon\varsigma_j$, since the supports are disjoint and $\varphi(0)=1$. By Lemmas~\ref{lem:embed} and~\ref{lem:bumps}, $\|g|_{B^e}\|_{\mathcal K_1(\cP^e)}\le2\varepsilon C_1\delta^{-s}=\varepsilon R/2=bBb_x/2$. Lemma~\ref{lem:finite} gives $N'\le N\le n$ elements $\sigma(\omega_i\cdot y+t_i)$ of $\cP^e$ and coefficients $a_i$ with $\sum_i|a_i|\le bBb_x/2$ and $\sum_ia_i\sigma(\omega_i\cdot y_j+t_i)=2\varepsilon\varsigma_j$ for every $j$.

For each $i$ let $w_i=(\omega_i,t_i)/\sqrt{1+t_i^2}\in S^{d-1}$. Then $w_i\cdot x_j=b_x(\omega_i\cdot y_j+t_i)/\sqrt{2(1+t_i^2)}$, so by positive homogeneity
\[
\sigma(\omega_i\cdot y_j+t_i)=\frac{\sqrt{2(1+t_i^2)}}{b_x}\,\sigma(w_i\cdot x_j).
\]
Put $\alpha_i=a_i\sqrt{2(1+t_i^2)}/b_x$, so that $\sum_i\alpha_i\sigma(w_i\cdot x_j)=2\varepsilon\varsigma_j$ and $\sum_i|\alpha_i|\le2\sum_i|a_i|/b_x\le bB$. Define $u\in\RR^n$ and $W\in\RR^{n\times d}$ by $u_i=\sgn(\alpha_i)\sqrt{|\alpha_i|b/B}$ and $i$th row $\sqrt{|\alpha_i|B/b}\,w_i$ for $i\le N'$, and $u_i=0$ with zero rows for $i>N'$. Then $|u|^2=(b/B)\sum_i|\alpha_i|\le b^2$, $\|W\|_F^2=(B/b)\sum_i|\alpha_i|\le B^2$, and $u_i\sigma(\sqrt{|\alpha_i|B/b}\,w_i\cdot x)=\alpha_i\sigma(w_i\cdot x)$. The function $f(x)=u^{T}\sigma(Wx)$ therefore lies in $\cF_n$ and satisfies $\varsigma_jf(x_j)=2\varepsilon\ge\varepsilon$ for every $j$.
\end{proof}

\section{A common bias}\label{sec:bias}

We prove Proposition~\ref{prop:bias} with the discretization of Section~\ref{sec:upper}.

\begin{proof}[Proof of Proposition~\ref{prop:bias}]
Let $f(x)=\sum_iu_i\sigma(w_i\cdot x-\beta)$ be in $\cF^\beta_n$. Since $0\le\sigma(z-\beta)\le|z|$, the argument of Lemma~\ref{lem:reduce} gives $|f|\le bBb_x$ on $X$ and $\fat_\varepsilon(\cF_n^\beta)=0$ when $R<1$. A neuron with $|w_i|\le\beta/b_x$ vanishes on $X$. For the others put $\omega_i=w_i/|w_i|$ and $t_i=-\beta/(|w_i|b_x)\in(-1,0]$; then for $y\in B^d$,
\[
u_i\sigma(w_i\cdot b_xy-\beta)=u_i|w_i|b_x\,\sigma(\omega_i\cdot y+t_i),
\]
so $f(b_xy)/(bBb_x)=\sum_ic_i\sigma(\omega_i\cdot y+t_i)$ with $\sum_i|c_i|\le1$ by the Cauchy--Schwarz inequality. Splitting the sum by the sign of $c_i$ writes this function as $P\int\sigma(\omega\cdot y+t)\,d\tau_+-Q\int\sigma(\omega\cdot y+t)\,d\tau_-$ with probability measures $\tau_\pm$ on $S^{d-1}\times[-1,1]$ and $P+Q\le1$. By~\cite[Theorem~2]{Siegel} with $k=1$ and $\alpha=0$ in its notation, each $\tau_\pm$ can be replaced by a probability measure with at most $m$ atoms at a uniform cost on $B^d$ of at most $A'_dm^{-1/2-3/(2d)}$. The functions $\sigma(\omega\cdot y+t)$ on $B^d$ satisfy $|\sigma(\omega\cdot y+t)-\sigma(\omega'\cdot y+t')|\le|\omega-\omega'|+|t-t'|$, so the discretization in the proof of Lemma~\ref{lem:Hcover}, with a net of $S^{d-1}\times[-1,1]$ in place of $N_1$ and no linear term, gives a uniform $\eta$-cover of the rescaled class with logarithm at most $C_d\eta^{-2d/(d+3)}\log(2/\eta)$. Lemma~\ref{lem:cover} with $\eta=1/(2R)$ completes the proof as in Section~\ref{sec:upper}.
\end{proof}

\section*{Acknowledgements and funding}
This work was done by multiple AI agents, Claude (Anthropic) and Codex (OpenAI), connected to Math Brain~\cite{MathBrainRepo} and working together inside the Brain framework~\cite{BrainFramework}, a shared, persistent memory in which the agents record claims, proofs, failed attempts and methods, retrieve one another's results, coordinate who works on what, and audit each other's work. For more on the details behind the Brain and Math Brain, see the Math Brain paper~\cite{MathBrainPaper}, the Math Brain repository~\cite{MathBrainRepo} and the Brain framework repository~\cite{BrainFramework}.

The research presented here was supported by the European Research Council (ERC) under the European Union's Horizon Europe research and innovation program (grant agreement No.~101041711) and by the Simons Foundation as part of the Collaboration on the Mathematical and Scientific Foundations of Deep Learning. It was also supported by Heights Labs, by Convex Nexus Capital, by the Israel Science Foundation (grant number 2258/19), and by the Israel Science Foundation (ISF Grant 4101/25).

\section*{Declaration of competing interest}
The author holds equity in Heights Labs and in Convex Nexus Capital and has a management role at Convex Nexus Capital. These entities are included in the funding declaration above.

\section*{Declaration of generative AI and AI-assisted technologies}
In November 2026 Claude Opus~5.5 (Anthropic), working in Claude Code and connected to Math Brain with the other agents named in the acknowledgements, developed the arguments, checked them against the cited theorems and wrote the paper. The question is that of Vardi, Shamir and Srebro~\cite{VSS}. The upper bound combines the connection between ReLU networks and zonoids due to Bach~\cite{Bach} with the approximation of zonoids by zonotopes of Matou\v{s}ek~\cite{Matousek} and Siegel~\cite{Siegel}. The lower bound rests on the Sobolev embedding of Mao, Siegel and Xu~\cite{MSX} and is new. The author directed the work and is responsible for the content of the paper and for decisions concerning its dissemination.

\section*{Data availability}
No data were generated. The \LaTeX{} source of the paper is available at \url{https://github.com/yspennstate/relu-fat-shattering-fixed-dimension/tree/v1.0}.

\appendix
\section{Notation}\label{app:notation}
For each symbol the table gives the place where it first appears. Letters used only inside one proof are marked as local to it.

{\small
\begin{longtable}{@{}p{0.2\textwidth}p{0.5\textwidth}l@{}}
\toprule
Symbol & Meaning & First appears \\ \midrule
\endhead
$d$ & input dimension, $d\ge2$ & Section~\ref{sec:intro} \\
$n$ & width, the number of rows of $W$ & Section~\ref{sec:intro} \\
$b$ & bound on the Euclidean norm of the output weights $u$ & Section~\ref{sec:intro} \\
$B$ & bound on the Frobenius norm of the inner weights $W$ & Section~\ref{sec:intro} \\
$b_x$ & radius of the input ball & Section~\ref{sec:intro} \\
$u$, $W$ & output weight vector and inner weight matrix of a network in $\cF_n$ & Section~\ref{sec:intro} \\
$\sigma$ & the ReLU function $\sigma(z)=\max\{z,0\}$ & Section~\ref{sec:intro} \\
$X$ & the ball $\{x\in\RR^d:|x|\le b_x\}$ & Section~\ref{sec:intro} \\
$\cF_n$ & networks $x\mapsto u^{T}\sigma(Wx)$ with $|u|\le b$ and $\|W\|_F\le B$ & Section~\ref{sec:intro} \\
$\varepsilon$ & scale of shattering & Section~\ref{sec:intro} \\
$N$ & number of shattered points & Section~\ref{sec:intro} \\
$x_j$ & the shattered points $x_1,\dots,x_N$ & Section~\ref{sec:intro} \\
$\hat g$ & Fourier transform of $g$ & Section~\ref{sec:prelim} \\
$W^s(L_2)$ & the Sobolev space of~\cite{MSX}, $W^s(L_2(\RR^e))$ and $W^s(L_2(B^e))$ & Section~\ref{sec:prelim} \\
$\tau'$, $\tau'_\pm$ & discrete measures with at most $m$ atoms & Lemma~\ref{lem:zonoid} \\
$h_{Z_m}$ & support function of the zonotope $Z_m$ & Lemma~\ref{lem:planar} \\
$C_e$ & constant of the embedding & Lemma~\ref{lem:embed} \\
$C_{e,\varphi}$ & constant in the Sobolev bound for sums of bumps & Lemma~\ref{lem:bumps} \\
$r$ & levels $r_1,\dots,r_N$ of a shattering & Section~\ref{sec:intro} \\
$s$ & the Sobolev order $(e+3)/2=(d+2)/2$ & Section~\ref{sec:intro} \\
$f$ & an element of a function class & Section~\ref{sec:intro} \\
$\fat$ & fat-shattering dimension, $\fat_\varepsilon$ at scale $\varepsilon$ & Section~\ref{sec:intro} \\
$R$ & the ratio $bBb_x/\varepsilon$ & Section~\ref{sec:intro} \\
$C$, $c$ & constants; a subscript $d$ or $e$ marks what they depend on & Section~\ref{sec:intro} \\
$w$ & a row $w_i$ of $W$, or a weight vector; in the proof of part (b) of the main theorem a unit vector $w_i$, the row being $\sqrt{|\alpha_i|B/b}\,w_i$ & Section~\ref{sec:intro} \\
$e$ & a dimension, the sphere $S^{e-1}\subset\RR^e$; $e=d-1$ in the lower bound & Section~\ref{sec:intro} \\
$\delta$ & bump radius in the lower bound & Section~\ref{sec:intro} \\
$\eta$ & scale of a uniform or $L_2$ cover & Section~\ref{sec:intro} \\
$\cF_n^{\mathrm{sp}}$ & as $\cF_n$, with the spectral norm $\|W\|\le B$ & Section~\ref{sec:intro} \\
$\cF_n^{\beta}$ & networks $x\mapsto u^{T}\sigma(Wx-\beta\mathbf 1)$ with a common bias $\beta\ge0$, where $\mathbf 1\in\RR^n$ is the vector of ones & Section~\ref{sec:intro} \\
$\beta$ & the common bias of the neurons of $\cF_n^{\beta}$ & Section~\ref{sec:intro} \\
$c_i$ & coefficients of a combination $\sum_ic_i\sigma(\omega_i\cdot x)$ of unit-norm neurons & Section~\ref{sec:intro} \\
$\omega$ & unit direction of a neuron, $\omega_i$ & Section~\ref{sec:intro} \\
$t$ & bias of a neuron, $t_i$ & Section~\ref{sec:intro} \\
$m$ & the width of the networks of Liu, Dadi and Cevher; the number of atoms of the discrete measure in Lemma~\ref{lem:zonoid} & Section~\ref{sec:intro} \\
$\tau$ & a finite positive measure on the sphere & Section~\ref{sec:upper} \\
$S^{d-1}$, $S^{e-1}$ & unit spheres of $\RR^d$ and $\RR^e$ & Section~\ref{sec:intro} \\
$\cP^e$ & neurons $\sigma(\omega\cdot y+t)$, $\omega\in S^{e-1}$, $t\in[-1,1]$, on the unit ball & Section~\ref{sec:intro} \\
$a$ & coefficients $a_i$ of a finite combination; in Section~\ref{sec:intro} also a candidate bound in the infimum after Proposition~\ref{prop:bias} & Section~\ref{sec:intro} \\
$g$ & an element of $L_2$ or of a cover & Section~\ref{sec:intro} \\
$B^e$ & closed unit ball of $\RR^e$ & Section~\ref{sec:intro} \\
$B_1(\cP^e)$ & closed absolutely convex hull of $\cP^e$ in $L_2(B^e)$ & Section~\ref{sec:prelim} \\
$\|\cdot\|_{\mathcal K_1(\cP^e)}$, $\mathcal K_1(\cP^e)$ & variation norm and variation space of ReLU networks with biases & Section~\ref{sec:intro} \\
$\xi$ & frequency variable of the Fourier transform & Section~\ref{sec:prelim} \\
$\|\cdot\|_{H^s}$ & Sobolev norm defined by the Fourier transform & Section~\ref{sec:prelim} \\
$G$ & a finite cover & Section~\ref{sec:intro} \\
$h_\tau$ & support function $\int|x\cdot y|\,d\tau(y)$ of a zonoid & Section~\ref{sec:upper} \\
$y$ & a point of the hyperplane in the lower bound, of the sphere or of the unit ball & Section~\ref{sec:intro} \\
$w'$ & the first $d-1$ coordinates of a weight vector $w=(w',w_d)$ & Section~\ref{sec:intro} \\
$\cF$ & a class of real functions & Section~\ref{sec:intro} \\
$A_d$ & constant in the zonoid approximation & Lemma~\ref{lem:zonoid} \\
$v$ & the vectors $v_i$ of a zonotope with $\sum_i|v_i|\le1$; unit vectors $v_i$ in the proof of the upper bound; the vector $(g(y_j))_j$ of values in the lower bound & Lemma~\ref{lem:planar} \\
$\mathcal H$ & normalized class $\sum_ic_i\sigma(\omega_i\cdot x)$ with $\sum_i|c_i|\le1$ & Lemma~\ref{lem:Hcover} \\
$P$, $Q$ & total positive and negative coefficient mass & Lemma~\ref{lem:Hcover} \\
$\mathbf m$ & linear part $\tfrac12\sum_ic_i\omega_i$ & Lemma~\ref{lem:Hcover} \\
$\Delta$ & mesh $\eta/32$ of the nets & Lemma~\ref{lem:Hcover} \\
$\gamma$ & the exponent $(d+2)/(2(d-1))$ & Lemma~\ref{lem:Hcover} \\
$\varphi$ & a smooth bump supported in the unit ball & Lemma~\ref{lem:bumps} \\
$\varsigma$ & sign vector of a shattering; the signs $\varsigma_j$ of the bumps in Section~\ref{sec:lower} & Section~\ref{sec:intro} \\
$M$ & a bound on the variation norm & Lemma~\ref{lem:finite} \\
$K$ & the compact set of neuron value vectors & Lemma~\ref{lem:finite} \\
$\lambda$ & least scaling with $v\in\lambda\operatorname{conv}K$ & Lemma~\ref{lem:finite} \\
$C'_d$ & the threshold for $R$ in part (b) of the main theorem & Section~\ref{sec:intro} \\
$C_1$ & a constant in the proof of part (b) of the main theorem & Section~\ref{sec:lower} \\
$A'_d$ & constant in the proof of Proposition~\ref{prop:bias} & Section~\ref{sec:upper} \\
$\alpha$ & a turning angle; rescaled outer coefficients $\alpha_i$ in the lower bound; a multi-index in Section~\ref{sec:lower}; in the proof of Proposition~\ref{prop:bias} the parameter $\alpha=0$ of~\cite[Theorem~2]{Siegel} & Lemma~\ref{lem:planar} \\
$Z$, $Z_m$, $\Lambda$, $p$, $q$, $\Gamma$, $\mu$, $\zeta$ & local to one proof & Lemma~\ref{lem:planar} \\
$N_1$, $N_2$, $N_3$ & nets, local to one proof & Lemma~\ref{lem:Hcover} \\
$k$ & an index of summation; the power of the ReLU$^k$ dictionary of Mao, Siegel and Xu; an integer Sobolev order in the lower bound & Section~\ref{sec:intro} \\
$h$ & a combination $\sum_ic_i\sigma(\omega_i\cdot x)$ in the class $\mathcal H$ & Section~\ref{sec:upper} \\
$\psi$, $\ell$ & local to one proof & Lemma~\ref{lem:planar} \\
$N'$, $M'$ & local to one proof & Lemma~\ref{lem:finite} \\
$\theta$ & an angle in the proof of Lemma~\ref{lem:planar}; the interpolation exponent in Lemma~\ref{lem:bumps} & Lemma~\ref{lem:planar} \\
\bottomrule
\end{longtable}}


\begin{thebibliography}{99}
\bibitem{ABCH} N.~Alon, S.~Ben-David, N.~Cesa-Bianchi and D.~Haussler, \emph{Scale-sensitive dimensions, uniform convergence, and learnability}, J. ACM 44 (1997), 615--631.
\bibitem{Bach} F.~Bach, \emph{Breaking the curse of dimensionality with convex neural networks}, J. Mach. Learn. Res. 18 (2017), 1--53, arXiv:1412.8690.
\bibitem{Bartlett98} P.~L.~Bartlett, \emph{The sample complexity of pattern classification with neural networks: the size of the weights is more important than the size of the network}, IEEE Trans. Inform. Theory 44 (1998), 525--536.
\bibitem{BFT} P.~L.~Bartlett, D.~J.~Foster and M.~Telgarsky, \emph{Spectrally-normalized margin bounds for neural networks}, Adv. Neural Inf. Process. Syst. 30 (2017), arXiv:1706.08498.
\bibitem{BLM} J.~Bourgain, J.~Lindenstrauss and V.~Milman, \emph{Approximation of zonoids by zonotopes}, Acta Math. 162 (1989), 73--141.
\bibitem{GRS} N.~Golowich, A.~Rakhlin and O.~Shamir, \emph{Size-independent sample complexity of neural networks}, Inf. Inference 9 (2020), 473--504, arXiv:1712.06541.
\bibitem{KS} M.~J.~Kearns and R.~E.~Schapire, \emph{Efficient distribution-free learning of probabilistic concepts}, J. Comput. System Sci. 48 (1994), 464--497.
\bibitem{LQWP} T.~Liang, D.~Qiao, Y.-X.~Wang and R.~Parhi, \emph{Stable minima of ReLU neural networks suffer from the curse of dimensionality: the neural shattering phenomenon}, preprint (2025), arXiv:2506.20779.
\bibitem{LDC} F.~Liu, L.~Dadi and V.~Cevher, \emph{Learning with norm constrained, over-parameterized, two-layer neural networks}, J. Mach. Learn. Res. 25 (2024), 1--42.
\bibitem{Makovoz} Y.~Makovoz, \emph{Random approximants and neural networks}, J. Approx. Theory 85 (1996), 98--109.
\bibitem{MSX} T.~Mao, J.~W.~Siegel and J.~Xu, \emph{Approximation rates for shallow ReLU$^k$ neural networks on Sobolev spaces via the Radon transform}, preprint (2025), arXiv:2408.10996v3.
\bibitem{MV} D.~E.~McClure and R.~A.~Vitale, \emph{Polygonal approximation of plane convex bodies}, J. Math. Anal. Appl. 51 (1975), 326--358.
\bibitem{Matousek} J.~Matou\v{s}ek, \emph{Improved upper bounds for approximation by zonotopes}, Acta Math. 177 (1996), 55--73.
\bibitem{NTS} B.~Neyshabur, R.~Tomioka and N.~Srebro, \emph{Norm-based capacity control in neural networks}, Proc. 28th Conf. Learning Theory, PMLR 40 (2015), 1376--1401, arXiv:1503.00036.
\bibitem{PN} R.~Parhi and R.~D.~Nowak, \emph{Near-minimax optimal estimation with shallow ReLU neural networks}, IEEE Trans. Inform. Theory 69 (2023), 1125--1140.
\bibitem{Siegel} J.~W.~Siegel, \emph{Optimal approximation of zonoids and uniform approximation by shallow neural networks}, preprint (2025), arXiv:2307.15285v3.
\bibitem{SX} J.~W.~Siegel and J.~Xu, \emph{Sharp bounds on the approximation rates, metric entropy, and $n$-widths of shallow neural networks}, Found. Comput. Math. 24 (2024), 481--537, arXiv:2101.12365.
\bibitem{VSS} G.~Vardi, O.~Shamir and N.~Srebro, \emph{The sample complexity of one-hidden-layer neural networks}, Adv. Neural Inf. Process. Syst. 35 (2022), arXiv:2202.06233.
\bibitem{MathBrainRepo} Y.~Shmalo, \emph{Math Brain: public snapshot, code and tests}, 2026, \url{https://github.com/yspennstate/math-brain-public/tree/8c9a6bd478c4ec4082347151b7ae16db50c917dd}.
\bibitem{MathBrainPaper} Y.~Shmalo, \emph{Math Brain: shared, auditable memory for multi-agent mathematical research}, October 2026, doi:10.5281/zenodo.24150837, \url{https://github.com/yspennstate/math-brain-public/tree/main/papers/math-brain-2026-10-24}.
\bibitem{BrainFramework} Y.~Shmalo, \emph{Brain Framework}, public software repository, 2026, \url{https://github.com/yspennstate/brain-framework-public/tree/5629e644cdd510173ae97e9aa92642dabd22dc57}.
\end{thebibliography}
\end{document}
